\documentclass[11pt,a4paper]{article}
\usepackage[margin=18mm]{geometry}
\usepackage{fontspec,microtype,graphicx,xcolor,tabularx,hyperref}
\setmainfont{Arial}
\setsansfont{Arial}
\setlength{\parindent}{0pt}
\setlength{\parskip}{5pt}
\pagestyle{empty}
\definecolor{accent}{HTML}{2563EB}
\definecolor{good}{HTML}{15803D}
% Generated from final CSV tables before release.
\newcommand{\MainResultShort}{KAN не ускоряет grokking универсально, а сдвигает границы между фазами обучения.}
\newcommand{\MainResultLong}{В полной матрице MLP дала 30/40 классических запусков, KAN --- 8/40. Локальный контраст возникает при $p=31$, 35\% train: MLP дала 0/10, KAN --- 8/10. При 50\% данных KAN обобщалась совместно или со слабой задержкой, а все 20 запусков $p=53$ получили класс weak delay.}
\newcommand{\KanFourierMem}{0.175}
\newcommand{\KanFourierGen}{0.828}
\newcommand{\KanFourierFinal}{0.952}
\newcommand{\MlpFourierMem}{0.075}
\newcommand{\MlpFourierFinal}{0.527}
\newcommand{\TopTenDelta}{-0.078}
\newcommand{\RandomTenDelta}{-0.047}
\newcommand{\BottomTenDelta}{0.000}
\newcommand{\RandomFourierFinal}{0.952}
\newcommand{\SymmetryFourierFinal}{0.165}
\newcommand{\PairedTableRows}{
KAN & 31 & 0.35 & 0.1 & 4/5 & 1 без обобщения & 4/5; 4.6k \\
KAN & 31 & 0.35 & 0.3 & 4/5 & 1 без обобщения & 4/5; 2.8k \\
KAN & 31 & 0.50 & 0.1 & 0/5 & weak 4, joint 1 & 5/5; 0.4k \\
KAN & 31 & 0.50 & 0.3 & 0/5 & weak 2, joint 3 & 5/5; 0.2k \\
KAN & 53 & 0.30 & 0.1 & 0/5 & weak 5 & 5/5; 3.9k \\
KAN & 53 & 0.30 & 0.3 & 0/5 & weak 5 & 5/5; 2.6k \\
KAN & 53 & 0.35 & 0.1 & 0/5 & weak 5 & 5/5; 2.3k \\
KAN & 53 & 0.35 & 0.3 & 0/5 & weak 5 & 5/5; 1.6k \\
MLP & 31 & 0.35 & 0.1 & 0/5 & 5 без обобщения & 0/5 \\
MLP & 31 & 0.35 & 0.3 & 0/5 & 5 без обобщения & 0/5 \\
MLP & 31 & 0.50 & 0.1 & 5/5 & -- & 5/5; 5.1k \\
MLP & 31 & 0.50 & 0.3 & 5/5 & -- & 5/5; 5.3k \\
MLP & 53 & 0.30 & 0.1 & 5/5 & -- & 5/5; 10.3k \\
MLP & 53 & 0.30 & 0.3 & 5/5 & -- & 5/5; 13.4k \\
MLP & 53 & 0.35 & 0.1 & 5/5 & -- & 5/5; 6.2k \\
MLP & 53 & 0.35 & 0.3 & 5/5 & -- & 5/5; 7.2k \\
}


\begin{document}
{\color{accent}\small КРАТКОЕ ОПИСАНИЕ ИССЛЕДОВАНИЯ}\par
{\LARGE\bfseries Grokking модульной арифметики: MLP против spline-KAN}\par
{\small Автор: Варфоломеев Матвей Денисович}\par
\vspace{3mm}

\textbf{Вопрос.} Как меняются запоминание и обобщение, если в задаче модульного
сложения заменить quadratic MLP на основе архитектуры Громова на spline-KAN
сопоставимого размера? Оптимизация адаптирована к сравнению на AdamW; работа
строит карту фаз, а не доказывает превосходство KAN.

\textbf{Протокол.} После exploratory-поиска на $p=97$ зафиксированы критерий и
матрица основной серии. Классический grokking требует устойчивого раннего
train-запоминания, низкого test в этот момент, плато не менее 1000 шагов и лишь
затем устойчивых 99\% test. Выполнены 80 запусков: четыре условия,
две архитектуры, два weight decay и пять seed; незавершившиеся результаты
цензурируются.

\textbf{Методы.} Модели получают one-hot $(a,b)$ и предсказывают
$(a+b)\bmod p$ при общих MSE, full-batch AdamW, \texttt{lr=0.001}, split и
бюджете. В точках запоминания, обобщения и финала измеряются
Fourier-диагональная мощность и PCA participation rank. Для пяти KAN зануляются
1\%, 5\% и 10\% top, random и bottom рёбер; random повторяется десять раз.

\textbf{Главный результат.} \MainResultLong

\textbf{Поддерживающие признаки.} У KAN медианная Fourier-концентрация растёт с
\KanFourierMem{} при запоминании до \KanFourierFinal{} в финале. Удаление 10\%
рёбер меняет test accuracy на \TopTenDelta{} для top, \RandomTenDelta{} для
random и \BottomTenDelta{} для bottom. Top сильнее bottom в 5/5 моделей, но
сильнее random-медианы лишь в 4/5.

\textbf{Проверка вывода.} $p=97$ исключён после exploratory-этапа: его
кривые не показывали убедительного grokking. Строгий критерий изменил класс
части 99\%-запусков. На symmetry-safe split KAN дала 0/5, а MLP в своём
режиме $p=31$, 50\% --- 5/5: результат KAN зависит от split.

\textbf{Воспроизводимость.} Репозиторий фиксирует критерий фаз, матрицу,
конфигурации, истории всех запусков и отрицательные результаты. Итоговые таблицы
строятся непосредственно из сохранённых метрик; контрольные суммы позволяют
проверить состав комплекта.

\textbf{Ограничения.} Пять seed дают широкие Wilson intervals;
Fourier-структура не доказывает единственный алгоритм; сплайн вне наблюдавшихся
one-hot значений не интерпретируется. Сохранены код, 80 историй и таблицы.

\vfill
{\small Основа: Gromov (2023), Power et al. (2022), Nanda et al. (2023),
Liu et al. (2024). Полный воспроизводимый протокол находится в репозитории.}
\end{document}
